\chapter{2003年天津卷（理）}
\section{单选题}
\begin{problem}
\(\frac{1 - \sqrt{3}\i}{(\sqrt{3} + \i)^2} =\)（\quad）

\choices
    {\(\frac{1}{4} + \frac{\sqrt{3}}{4}\i\)}
    {\(-\frac{1}{4} - \frac{\sqrt{3}}{4}\i\)}
    {\(\frac{1}{2} + \frac{\sqrt{3}}{2}\i\)}
    {\(-\frac{1}{2} - \frac{\sqrt{3}}{2}\i\)}
\end{problem}

\begin{answer}
B
\end{answer}

\begin{solution}
先计算分母，
\[
(\sqrt{3}+\i)^2=2+2\sqrt{3}\i
\]
于是
\[
\frac{1-\sqrt{3}\i}{(\sqrt{3}+\i)^2}=\frac{1-\sqrt{3}\i}{2(1+\sqrt{3}\i)}=\frac{(1-\sqrt{3}\i)^2}{2(1+3)}=-\frac14-\frac{\sqrt{3}}4\i
\]
故选 B.
\end{solution}

\begin{problem}
已知 \(x \in \left(-\frac{\pi}{2}, 0\right)\)，\(\cos x = \frac{4}{5}\)，则 \(\tan 2x =\)（\quad）

\choices
    {\(\frac{7}{24}\)}
    {\(-\frac{7}{24}\)}
    {\(\frac{24}{7}\)}
    {\(-\frac{24}{7}\)}
\end{problem}

\begin{answer}
D
\end{answer}

\begin{solution}
因为 \(x\in\left(-\frac\pi2,0\right)\)，所以 \(\sin x<0\)，从而
\(\sin x=-\sqrt{1-\cos^2x}=-\frac35\)，\(\tan x=-\frac34\).
因此
\[
\tan2x=\frac{2\tan x}{1-\tan^2x}=\frac{-\frac32}{1-\frac9{16}}=-\frac{24}{7}
\]
故选 D.
\end{solution}

\begin{problem}
设函数 \( f(x) = \begin{cases}
2^{-x} - 1, & x \leqslant 0,\\
x^{\frac{1}{2}}, & x > 0,
\end{cases} \) 若 \( f(x_0) > 1 \)，则 \( x_0 \) 的取值范围是（\quad）

\choices
    {\( (-1,1) \)}
    {\((-1, +\infty)\)}
    {\((- \infty, -2) \cup (0, +\infty)\)}
    {\((- \infty, -1) \cup (1, +\infty)\)}
\end{problem}

\begin{answer}
D
\end{answer}

\begin{solution}
当 \(x_0\leqslant0\) 时，
\[
f(x_0)>1\Longleftrightarrow 2^{-x_0}-1>1\Longleftrightarrow x_0<-1
\]
当 \(x_0>0\) 时，
\[
f(x_0)>1\Longleftrightarrow \sqrt{x_0}>1\Longleftrightarrow x_0>1
\]
所以
\[
x_0\in(-\infty,-1)\cup(1,+\infty)
\]
故选 D.
\end{solution}

\begin{problem}
\(O\) 是平面上一定点，\(A\)，\(B\)，\(C\) 是平面上不共线的三个点，动点 \(P\) 满足 \(\overrightarrow{OP} = \overrightarrow{OA} + \lambda \left(\frac{\overrightarrow{AB}}{|\overrightarrow{AB}|} + \frac{\overrightarrow{AC}}{|\overrightarrow{AC}|}\right)\)（\(\lambda \in [0, +\infty)\)），则 \(P\) 的轨迹一定通过 \(\triangle ABC\) 的（\quad）

\choices
    {外心}
    {内心}
    {重心}
    {垂心}
\end{problem}

\begin{answer}
B
\end{answer}

\begin{solution}
设 \(a=|BC|\)，\(b=|CA|\)，\(c=|AB|\)，\(I\) 为三角形 \(ABC\) 的内心. 由内心的重心式表示，
\[
\overrightarrow{AI}=\frac{b\overrightarrow{AB}+c\overrightarrow{AC}}{a+b+c}=\frac{bc}{a+b+c}\left(\frac{\overrightarrow{AB}}{c}+\frac{\overrightarrow{AC}}{b}\right)
\]
取 \(\lambda=\dfrac{bc}{a+b+c}>0\)，则题设中的点 \(P\) 为 \(I\). 故轨迹一定通过内心，选 B.
\end{solution}

\begin{problem}
函数 \(y = \ln \frac{x + 1}{x - 1}\)（\(x \in (1, +\infty)\)）的反函数为（\quad）

\choices
    {\(y = \frac{\e^x - 1}{\e^x + 1}\)（\(x \in (0, +\infty)\)）}
    {\(y = \frac{\e^x + 1}{\e^x - 1}\)（\(x \in (0, +\infty)\)）}
    {\(y = \frac{\e^x - 1}{\e^x + 1}\)（\(x \in (-\infty, 0)\)）}
    {\(y = \frac{\e^x + 1}{\e^x - 1}\)（\(x \in (-\infty, 0)\)）}
\end{problem}

\begin{answer}
B
\end{answer}

\begin{solution}
令原函数中的自变量为 \(t\)，函数值为 \(y\)，则
\[
e^y=\frac{t+1}{t-1}\Longrightarrow e^y(t-1)=t+1\Longrightarrow t=\frac{e^y+1}{e^y-1}
\]
当 \(t>1\) 时，原函数的值域为 \((0,+\infty)\). 交换自变量与函数值，反函数为
\(y=\frac{e^x+1}{e^x-1}\)，\(x\in(0,+\infty)\).
故选 B.
\end{solution}

\begin{problem}
棱长为 \(a\) 的正方体中，连结相邻面的中心，以这些线段为棱的八面体的体积为（\quad）

\choices
    {\(\frac{a^3}{3}\)}
    {\(\frac{a^3}{4}\)}
    {\(\frac{a^3}{6}\)}
    {\(\frac{a^3}{12}\)}
\end{problem}

\begin{answer}
C
\end{answer}

\begin{solution}
以正方体中心为原点，三条棱方向为坐标轴，则六个相邻面的中心为
\(\left(\pm\frac a2,0,0\right)\)，\(\left(0,\pm\frac a2,0\right)\)，\(\left(0,0,\pm\frac a2\right)\).
八面体在第一卦限内的部分是截距均为 \(a/2\) 的三棱锥，体积为 \(\dfrac16\left(\dfrac a2\right)^3\). 八个卦限全等，所以八面体体积为
\[
8\cdot\frac16\left(\frac a2\right)^3=\frac{a^3}{6}
\]
故选 C.
\end{solution}

\begin{problem}
设 \(a > 0\)，\(f(x) = ax^2 + bx + c\)，曲线 \(y = f(x)\) 在点 \(P(x_0, f(x_0))\) 处切线的倾斜角的取值范围为 \(\left[0, \frac{\pi}{4}\right]\)，则 \(P\) 到曲线 \(y = f(x)\) 对称轴距离的取值范围为（\quad）

\choices
    {\(\left[0, \frac{1}{a}\right]\)}
    {\(\left[0, \frac{1}{2a}\right]\)}
    {\(\left[0, \left|\frac{b}{2a}\right|\right]\)}
    {\(\left[0, \left|\frac{b-1}{2a}\right|\right]\)}
\end{problem}

\begin{answer}
B
\end{answer}

\begin{solution}
抛物线在 \(P(x_0,f(x_0))\) 处的切线斜率为
\[
f'(x_0)=2ax_0+b
\]
切线倾斜角属于 \(\left[0,\dfrac\pi4\right]\)，所以
\[
0\leqslant2ax_0+b\leqslant1
\]
由 \(a>0\) 得
\[
-\frac b{2a}\leqslant x_0\leqslant\frac{1-b}{2a}
\]
抛物线对称轴为直线 \(x=-\dfrac b{2a}\)，点 \(P\) 到该直线的距离为 \(x_0+\dfrac b{2a}\)，故距离的取值范围为
\[
\left[0,\frac1{2a}\right]
\]
故选 B.
\end{solution}

\begin{problem}
已知方程 \((x^{2} - 2x + m)(x^{2} - 2x + n) = 0\) 的四个根组成一个首项为 \(\frac{1}{4}\) 的等差数列，则 \(|m - n| =\)（\quad）

\choices
    {1}
    {\(\frac{3}{4}\)}
    {\(\frac{1}{2}\)}
    {\(\frac{3}{8}\)}
\end{problem}

\begin{answer}
C
\end{answer}

\begin{solution}
设等差数列四项依次为
\(r_1=\frac14\)，\(r_2=\frac14+d\)，\(r_3=\frac14+2d\)，\(r_4=\frac14+3d\).
两个二次因式的根分别成对出现，且每一对根的和都等于 \(2\). 四项的三种配对中，只有 \((r_1,r_4)\)、\((r_2,r_3)\) 能使两对的和相等；于是
\(r_1+r_4=\frac12+3d=2\)，\(d=\frac12\).
因此四个根为 \(\dfrac14\)，\(\dfrac34\)，\(\dfrac54\)，\(\dfrac74\)，两个二次因式的常数项分别为
\(m=\frac14\cdot\frac74=\frac7{16}\)，\(n=\frac34\cdot\frac54=\frac{15}{16}\)
二者次序可交换，故 \(|m-n|=\dfrac12\)，选 C.
\end{solution}

\begin{problem}
已知双曲线中心在原点且一个焦点为 \(F(\sqrt{7}, 0)\)，直线 \(y = x - 1\) 与其相交于 \(M\)，\(N\) 两点，\(MN\) 中点的横坐标为 \(-\frac{2}{3}\)，则此双曲线的方程是（\quad）

\choices
    {\(\frac{x^2}{3} -\frac{y^2}{4} = 1\)}
    {\(\frac{x^2}{4} -\frac{y^2}{3} = 1\)}
    {\(\frac{x^2}{5} -\frac{y^2}{2} = 1\)}
    {\(\frac{x^2}{2} -\frac{y^2}{5} = 1\)}
\end{problem}

\begin{answer}
D
\end{answer}

\begin{solution}
设双曲线方程为 \(\dfrac{x^2}{a^2}-\dfrac{y^2}{b^2}=1\)，则由焦点为 \((\sqrt7,0)\) 得
\[
a^2+b^2=7
\]
将直线 \(y=x-1\) 代入双曲线方程，得关于 \(x\) 的二次方程
\[
\left(\frac1{a^2}-\frac1{b^2}\right)x^2+\frac2{b^2}x-\frac1{b^2}-1=0
\]
其两根为交点 \(M\)，\(N\) 的横坐标，故由韦达定理，弦中点横坐标为
\[
\frac12\cdot\frac{-\frac2{b^2}}{\frac1{a^2}-\frac1{b^2}}=\frac{a^2}{a^2-b^2}=-\frac23
\]
从而 \(5a^2=2b^2\). 结合 \(a^2+b^2=7\)，得 \(a^2=2\)，\(b^2=5\). 所以双曲线方程为
\[
\frac{x^2}{2}-\frac{y^2}{5}=1
\]
故选 D.
\end{solution}

\begin{problem}
已知长方形的四个顶点 \(A(0,0)\)，\(B(2,0)\)，\(C(2,1)\) 和 \(D(0,1)\)，一质点从 \(AB\) 的中点 \(P_0\) 沿与 \(AB\) 的夹角 \(\theta\) 的方向射到 \(BC\) 上的点 \(P_1\) 后，依次反射到 \(CD\)，\(DA\) 和 \(AB\) 上的点 \(P_2\)，\(P_3\) 和 \(P_4\)（入射角等于反射角）. 设 \(P_4\) 的坐标为 \((x_4,0)\)，若 \(1 < x_4 < 2\)，则 \(\tan \theta\) 的取值范围是（\quad）

\choices
    {\(\left(\frac{1}{3}, 1\right)\)}
    {\(\left(\frac{1}{3}, \frac{2}{3}\right)\)}
    {\( \left(\frac{2}{5},\frac{1}{2}\right) \)}
    {\( \left(\frac{2}{5},\frac{2}{3}\right) \)}
\end{problem}

\begin{answer}
C
\end{answer}

\begin{solution}
设 \(t=\tan\theta\). 由运动方向和连续反射过程可知 \(t>0\). 从 \(P_0=(1,0)\) 出发，第一次到达右边界时的竖直位移为 \(t\)，故
\[
P_1=(2,t)
\]
在 \(P_1\) 处反射后，水平方向反向而竖直方向不变，所以从 \(P_1\) 到 \(P_2\) 的斜率为 \(-t\). 令 \(P_2=(x_2,1)\)，则
\[
1-t=-t(x_2-2)\Longrightarrow x_2=3-\frac1t
\]
即 \(P_2=\left(3-\frac1t,1\right)\). 在 \(P_2\) 处反射后，竖直方向反向而水平方向不变；令 \(P_3=(0,y_3)\)，则从 \(P_2\) 到 \(P_3\) 的斜率为 \(t\)，从而
\[
y_3-1=t(0-x_2)\Longrightarrow y_3=1-t\left(3-\frac1t\right)=2-3t
\]
因此 \(P_3=(0,2-3t)\). 在 \(P_3\) 处再次反射后，令 \(P_4=(x_4,0)\)，由斜率为 \(-t\) 得
\[
-y_3=-tx_4\Longrightarrow x_4=\frac{y_3}{t}=\frac2t-3
\]
所以 \(P_4=\left(\frac2t-3,0\right)\).
题设给出 \(1<x_4<2\)，所以
\[
1<\frac2t-3<2\Longleftrightarrow \frac25<t<\frac12
\]
在此范围内，\(P_1\)，\(P_2\)，\(P_3\) 分别位于 \(BC\)，\(CD\)，\(DA\) 的内部，反射顺序和边界条件均满足. 因此 \(\tan\theta\) 的取值范围为 \(\left(\dfrac25,\dfrac12\right)\)，故选 C.
\end{solution}

\begin{problem}
\(\displaystyle\lim_{n\to \infty}\frac{\mathrm C_2^2 + \mathrm C_3^2 + \mathrm C_4^2 + \cdots + \mathrm C_n^2}{n(\mathrm C_2^1 + \mathrm C_3^1 + \mathrm C_4^1 + \cdots + \mathrm C_n^1)} =\)（\quad）

\choices
    {3}
    {\( \frac{1}{3} \)}
    {\( \frac{1}{6} \)}
    {6}
\end{problem}

\begin{answer}
B
\end{answer}

\begin{solution}
由组合数恒等式
\(\mathrm C_2^2+\mathrm C_3^2+\cdots+\mathrm C_n^2=\mathrm C_{n+1}^3\)，\(\mathrm C_2^1+\mathrm C_3^1+\cdots+\mathrm C_n^1=\frac{n(n+1)}2-1\)
原极限等于
\[
\lim_{n\to\infty}\frac{\frac{n(n+1)(n-1)}6}{n\left(\frac{n(n+1)}2-1\right)}=\lim_{n\to\infty}\frac{n^2-1}{3n^2+3n-6}=\frac13
\]
故选 B.
\end{solution}

\begin{problem}
一个四面体的所有棱长都为 \(\sqrt{2}\)，四个顶点在同一球面上，则此球的表面积为（\quad）

\choices
    {\(3\pi\)}
    {\(4\pi\)}
    {\(3\sqrt{3}\pi\)}
    {\(6\pi\)}
\end{problem}

\begin{answer}
A
\end{answer}

\begin{solution}
取正四面体的四个顶点为
\(\left(\frac12,\frac12,\frac12\right)\)，\(\left(\frac12,-\frac12,-\frac12\right)\)，\(\left(-\frac12,\frac12,-\frac12\right)\)，\(\left(-\frac12,-\frac12,\frac12\right)\).
任意两点相差的两个坐标各为 \(1\) 或 \(-1\)，所以棱长均为 \(\sqrt2\). 四点到原点的距离均为
\[
R=\sqrt{\frac14+\frac14+\frac14}=\frac{\sqrt3}{2}
\]
故原点为外接球球心，球的表面积为
\[
4\pi R^2=4\pi\cdot\frac34=3\pi
\]
故选 A.
\end{solution}

\section{填空题}
\begin{problem}
\(\left(x^{2} - \frac{1}{2x}\right)^{9}\) 的展开式中 \(x^{9}\) 系数是\fillinblank{}.
\end{problem}

\begin{answer}
\(-\frac{21}{2}\)
\end{answer}

\begin{solution}
展开式的通项为
\[
\mathrm C_9^k(x^2)^{9-k}\left(-\frac1{2x}\right)^k=\mathrm C_9^k\frac{(-1)^k}{2^k}x^{18-3k}
\]
要得到 \(x^9\) 项，必须有 \(18-3k=9\)，即 \(k=3\). 所以所求系数为
\[
\mathrm C_9^3\frac{(-1)^3}{2^3}=-\frac{84}{8}=-\frac{21}{2}
\]
\end{solution}

\begin{problem}
某公司生产三种型号的轿车，产量分别为 \(1200\text{ 辆}\)，\(6000\text{ 辆}\)和 \(2000\text{ 辆}\). 为检验该公司的产品质量，现用分层抽样的方法抽取 \(46\text{ 辆}\)进行检验，这三种型号的轿车依次应抽取\fillinblank{}，\fillinblank{}，\fillinblank{} \(\text{辆}\).
\end{problem}

\begin{answer}
\(6\)，\(30\)，\(10\)
\end{answer}

\begin{solution}
三种型号的轿车总数为 \(1200+6000+2000=9200\)，抽样比为 \(46/9200=1/200\). 所以三层分别抽取
\(1200\cdot\frac1{200}=6\)，\(6000\cdot\frac1{200}=30\)，\(2000\cdot\frac1{200}=10\).
故依次抽取 \(6\)，\(30\)，\(10\) 辆.
\end{solution}

\begin{problem}
某城市在中心广场建造一个花圃，花圃分为 \(6\text{ 个}\)部分（如图）. 现要栽种 \(4\text{ 种}\)不同颜色的花，每部分栽种一种且相邻部分不能栽种同样颜色的花，不同的栽种方法有\fillinblank{} 种.（以数字作答）

\bitmapfigure[max width=6.0cm]{img/2003/tianjin_science/q15_fig1.png}
\end{problem}

\begin{answer}
\(120\)
\end{answer}

\begin{solution}
区域 1 与其余五个区域都相邻，区域 2，3，4，5，6 依次还构成一个五边形环. 先给区域 1 选颜色，有 \(4\) 种选法；其余五个区域只能使用剩下的 \(3\) 种颜色.

固定区域 2 的颜色为 \(A\). 区域 3 有两种颜色可选. 若区域 3 的颜色为 \(B\)，则区域 4 取 \(A\) 时，区域 5、6 有 \(2\) 种排法；区域 4 取第三种颜色时有 \(3\) 种排法，所以固定区域 3 后有 \(5\) 种. 区域 3 有两种颜色，故固定区域 2 后有 \(2\times5=10\) 种；区域 2 再有 \(3\) 种颜色可选，所以外环共有 \(30\) 种.

因此总的栽种方法数为
\[
4\times30=120
\]
\end{solution}

\begin{problem}
下列 \(5\text{ 个}\)正方体图形中，\(l\) 是正方体的一条对角线，点 \(M\)，\(N\)，\(P\) 分别为其所在棱的中点，能得出 \(l \perp\) 面 \(MNP\) 的图形的序号是\fillinblank{}.（写出所有符合要求的图形序号）

\FigureLayoutDeclare{img/2003/tianjin_science/q16_fig1.png}{16.0cm}{20bp 37bp 57bp 4bp}
\bitmapfigure[max width=0.82\linewidth]{img/2003/tianjin_science/q16_fig1.png}
\end{problem}

\begin{answer}
\(1\)，\(4\)，\(5\)
\end{answer}

\begin{solution}
取正方体一个顶点为原点，三条棱为坐标轴，棱长取 1，并按图令对角线 \(l\) 的方向向量为 \(\bs d=(1,-1,-1)\). 对每个图形取出三个中点，计算平面 \(MNP\) 的法向量，可得下表；法向量允许同时乘以非零常数.

\begin{center}
\begin{tabular}{c|c|c|c|c}
图形 & \(M\) & \(N\) & \(P\) & \(MNP\) 的法向量 \\
\hline
1 & \(\left(0,1,\frac12\right)\) & \(\left(0,\frac12,1\right)\) & \(\left(\frac12,1,1\right)\) & \((1,-1,-1)\) \\
2 & \(\left(0,0,\frac12\right)\) & \(\left(1,0,\frac12\right)\) & \(\left(\frac12,0,0\right)\) & \((0,1,0)\) \\
3 & \(\left(0,\frac12,0\right)\) & \(\left(1,0,\frac12\right)\) & \(\left(\frac12,1,1\right)\) & \((1,1,-1)\) \\
4 & \(\left(0,\frac12,0\right)\) & \(\left(1,\frac12,1\right)\) & \(\left(0,0,\frac12\right)\) & \((1,-1,-1)\) \\
5 & \(\left(0,0,\frac12\right)\) & \(\left(1,\frac12,1\right)\) & \(\left(\frac12,1,0\right)\) & \((1,-1,-1)\)
\end{tabular}
\end{center}

直线垂直于平面当且仅当其方向向量平行于平面的法向量，因此符合要求的图形序号为 \(1\)，\(4\)，\(5\).
\end{solution}

\section{解答题}
\begin{problem}
已知函数 \( f(x) = 2\sin x (\sin x + \cos x) \).

\begin{enumerate}
\item 求函数 \( f(x) \) 的最小正周期和最大值；
\item 在给出的直角坐标系中，画出函数 \( y = f(x) \) 在区间 \(\left[-\frac{\pi}{2}, \frac{\pi}{2}\right]\) 上的图象.
\end{enumerate}
\bitmapfigure[max width=7.0cm]{img/2003/tianjin_science/q17_fig1.png}
\end{problem}

\begin{solution}
\begin{enumerate}
\item
\[
f(x)=2\sin^2x+2\sin x\cos x=1-\cos2x+\sin2x=1+\sqrt2\sin\left(2x-\frac\pi4\right)
\]
由于正弦函数的最小正周期为 \(2\pi\)，所以 \(f(x)\) 的最小正周期为 \(\pi\). 又因为 \(-1\leqslant\sin\left(2x-\dfrac\pi4\right)\leqslant1\)，故
\[
f_{\max}=1+\sqrt2
\]

\item 在区间 \(\left[-\dfrac\pi2,\dfrac\pi2\right]\) 内，函数图象经过
\(\left(-\frac\pi2,2\right)\)，\(\left(-\frac\pi4,0\right)\)，\(\left(-\frac\pi8,1-\sqrt2\right)\)，\((0,0)\)
\(\left(\frac\pi8,1\right)\)，\(\left(\frac{3\pi}8,1+\sqrt2\right)\)，\(\left(\frac\pi2,2\right)\).
其中 \(f'(x)=2\sin2x+2\cos2x\)，所以函数在 \(\left[-\dfrac\pi2,-\dfrac\pi8\right]\) 上递减，在 \(\left[-\dfrac\pi8,\dfrac{3\pi}8\right]\) 上递增，在 \(\left[\dfrac{3\pi}8,\dfrac\pi2\right]\) 上递减. 将上述各点按此单调性用光滑曲线连接，即得给定区间上的图象.
\end{enumerate}
\end{solution}

\begin{problem}
如图，在直三棱柱 \(ABC - A_{1}B_{1}C_{1}\) 中，底面是等腰直角三角形，\(\angle ACB = 90^{\circ}\)，侧棱 \(AA_{1} = 2\)，\(D\)，\(E\) 分别是 \(CC_{1}\) 与 \(A_{1}B\) 的中点，点 \(E\) 在平面 \(ABD\) 上的射影是 \(\triangle ABD\) 的重心 \(G\).

\begin{enumerate}
\item 求 \(A_{1}B\) 与平面 \(ABD\) 所成角的大小（结果用反三角函数值表示）；
\item 求点 \(A_{1}\) 到平面 \(AED\) 的距离.
\end{enumerate}
\bitmapfigure[max width=6.0cm]{img/2003/tianjin_science/q18_fig1.png}
\end{problem}

\begin{solution}
\begin{enumerate}
\item 设 \(AC=BC=s\)，建立空间直角坐标系，取
\(C=(0,0,0)\)，\(A=(s,0,0)\)，\(B=(0,s,0)\)，\(A_1=(s,0,2)\)，\(D=(0,0,1)\).
因为 \(E\) 是 \(A_1B\) 的中点，所以 \(E=\left(\dfrac s2,\dfrac s2,1\right)\). 三角形 \(ABD\) 的重心为
\[
G=\left(\frac s3,\frac s3,\frac13\right)
\]
由题意，\(\overrightarrow{EG}\) 垂直于平面 \(ABD\). 取平面内两个方向向量
\(\overrightarrow{AB}=(-s,s,0)\)，\(\overrightarrow{AD}=(-s,0,1)\).
其中 \(\overrightarrow{EG}=\left(-\dfrac s6,-\dfrac s6,-\dfrac23\right)\)，且它与 \(\overrightarrow{AB}\) 的内积恒为 0；再由
\[
\overrightarrow{EG}\cdot\overrightarrow{AD}=\frac{s^2}{6}-\frac23=0
\]
得 \(s=2\). 于是
\(A=(2,0,0)\)，\(B=(0,2,0)\)，\(D=(0,0,1)\)，\(A_1=(2,0,2)\).
平面 \(ABD\) 的一个法向量为
\[
\overrightarrow{AB}\times\overrightarrow{AD}=(2,2,4)
\]
而 \(\overrightarrow{A_1B}=(-2,2,-2)\)，所以直线与平面的夹角 \(\alpha\) 满足
\[
\sin\alpha=\frac{|\overrightarrow{A_1B}\cdot(\overrightarrow{AB}\times\overrightarrow{AD})|}{|\overrightarrow{A_1B}|\,|\overrightarrow{AB}\times\overrightarrow{AD}|}=\frac{\sqrt2}{3}
\]
故 \(\alpha=\arcsin\dfrac{\sqrt2}{3}\).

\item 由 \(A=(2,0,0)\)，\(E=(1,1,1)\)，\(D=(0,0,1)\)，平面 \(AED\) 的一个法向量为
\[
\overrightarrow{AE}\times\overrightarrow{AD}=(1,-1,2)
\]
故其方程为 \(x-y+2z-2=0\). 点 \(A_1=(2,0,2)\) 到此平面的距离为
\[
\frac{|2-0+2\times2-2|}{\sqrt{1^2+(-1)^2+2^2}}=\frac4{\sqrt6}=\frac{2\sqrt6}{3}
\]
\end{enumerate}
\end{solution}

\begin{problem}
设 \(a > 0\)，求函数 \(f(x) = \sqrt{x} - \ln (x + a)\)（\(x \in (0, +\infty)\)）的单调区间.
\end{problem}

\begin{solution}
\[
f'(x)=\frac1{2\sqrt{x}}-\frac1{x+a}=\frac{x+a-2\sqrt{x}}{2\sqrt{x}(x+a)}
\]
分母在 \(x>0\) 时恒为正. 令 \(t=\sqrt{x}>0\)，则分子为
\[
t^2-2t+a=(t-1)^2+a-1
\]
当 \(a\geqslant1\) 时，该式非负，且除至多一个点外为正，所以函数在 \((0,+\infty)\) 上单调递增.

当 \(0<a<1\) 时，分子在
\(t=1-\sqrt{1-a}\)，\(t=1+\sqrt{1-a}\)
处为零，并在两根外为正、两根之间为负. 将两根平方，得到
\(x_1=2-a-2\sqrt{1-a}\)，\(x_2=2-a+2\sqrt{1-a}\).
因此函数在 \((0,x_1)\) 和 \((x_2,+\infty)\) 上单调递增，在 \((x_1,x_2)\) 上单调递减.
\end{solution}

\begin{problem}
\(A\)，\(B\) 两个代表队进行乒乓球对抗赛，每队三名队员，\(A\) 队队员是 \(A_1\)，\(A_2\)，\(A_3\)，\(B\) 队队员是 \(B_1\)，\(B_2\)，\(B_3\)，按以往多次比赛的统计，对阵队员之间胜负概率如下：

\begin{center}
\begin{tabular}{|c|c|c|}
\hline
对阵队员 & \(A\) 队队员胜的概率 & \(A\) 队队员负的概率 \\
\hline
\(A_1\) 对 \(B_1\) & \(\frac{2}{3}\) & \(\frac{1}{3}\) \\
\hline
\(A_2\) 对 \(B_2\) & \(\frac{2}{5}\) & \(\frac{3}{5}\) \\
\hline
\(A_3\) 对 \(B_3\) & \(\frac{2}{5}\) & \(\frac{3}{5}\) \\
\hline
\end{tabular}
\end{center}

现按表中对阵方式出场，每场胜队得 \(1\text{ 分}\)，负队得 \(0\text{ 分}\)，设 \(A\) 队，\(B\) 队最后所得总分分别为 \(\xi\)，\(\eta\).

\begin{enumerate}
\item 求 \(\xi\)，\(\eta\) 的概率分布；
\item 求 \(E(\xi)\)，\(E(\eta)\).
\end{enumerate}
\end{problem}

\begin{solution}
\begin{enumerate}
\item 设三场比赛中 A 队获胜的指示变量之和为 \(\xi\)，并用 \(z\) 记录 A 队每赢一场所增加的分数. 按题意各场结果相互独立，故 \(\xi\) 的概率生成式为
\[
\left(\frac13+\frac23z\right)\left(\frac35+\frac25z\right)^2
\]
展开得
\[
\left(\frac13+\frac23z\right)\left(\frac9{25}+\frac{12}{25}z+\frac4{25}z^2\right)=\frac3{25}+\frac25z+\frac{28}{75}z^2+\frac8{75}z^3
\]
所以
\(P(\xi=0)=\frac3{25}\)，\(P(\xi=1)=\frac25\)，\(P(\xi=2)=\frac{28}{75}\)，\(P(\xi=3)=\frac8{75}\).
每场比赛恰有一队获胜，故 \(\eta=3-\xi\)，从而
\(P(\eta=0)=\frac8{75}\)，\(P(\eta=1)=\frac{28}{75}\)，\(P(\eta=2)=\frac25\)，\(P(\eta=3)=\frac3{25}\).

\item 由独立比赛中得分的期望可加，
\[
E(\xi)=\frac23+\frac25+\frac25=\frac{22}{15}
\]
又因为 \(\eta=3-\xi\)，所以
\[
E(\eta)=3-E(\xi)=\frac{23}{15}
\]
\end{enumerate}
\end{solution}

\begin{problem}
已知常数 \(a > 0\)，向量 \(\bs{c} = (0, a)\)，\(\bs{i} = (1, 0)\). 经过原点 \(O\) 以 \(\bs{c} + \lambda \bs{i}\) 为方向向量的直线与经过定点 \(A(0, a)\) 以 \(\bs{i} - 2\lambda \bs{c}\) 为方向向量的直线相交于 \(P\)，其中 \(\lambda \in \R\). 试问：是否存在两个定点 \(E\)，\(F\)，使得 \(|PE| + |PF|\) 为定值. 若存在，求出 \(E\)，\(F\) 的坐标；若不存在，说明理由.
\end{problem}

\begin{solution}
设第一条直线上的交点参数为 \(s\)，第二条直线上的交点参数为 \(t\)，则
\[
P=s(\lambda,a)=(t,a-2a\lambda t)
\]
因此 \(x=t=s\lambda\)，\(y=as\)，从而
\(y=a-2a\lambda x=a-2\lambda^2y\)，\(y=\frac{a}{1+2\lambda^2}\)，\(x=\frac{\lambda}{1+2\lambda^2}\).
消去 \(\lambda\)，得到
\[
2a^2x^2+y^2-ay=0
\]
即
\[
\frac{x^2}{1/8}+\frac{\left(y-\frac a2\right)^2}{a^2/4}=1
\]
当 \(\lambda\) 遍历 \(\R\) 时，\(P\) 遍历这个椭圆除去原点的部分；缺少一个极限点不影响椭圆焦点距离和的结论.

当 \(a=\dfrac{\sqrt2}{2}\) 时，椭圆退化为以 \((0,\dfrac a2)\) 为圆心的圆. 圆只有在两个焦点重合于圆心时才有相应的焦点距离和；题意要求两个定点，故此时不存在符合要求的两个定点.

当 \(0<a<\dfrac{\sqrt2}{2}\) 时，横轴为长轴，半长轴为 \(\dfrac{\sqrt2}{4}\)，焦距的一半为
\[
c=\sqrt{\left(\frac{\sqrt2}{4}\right)^2-\left(\frac a2\right)^2}=\frac12\sqrt{\frac12-a^2}
\]
所以可取
\(E=\left(\frac12\sqrt{\frac12-a^2},\frac a2\right)\)，\(F=\left(-\frac12\sqrt{\frac12-a^2},\frac a2\right)\)
且 \(|PE|+|PF|=\dfrac{\sqrt2}{2}\).

当 \(a>\dfrac{\sqrt2}{2}\) 时，纵轴为长轴，半长轴为 \(\dfrac a2\)，焦距的一半为
\[
c=\sqrt{\left(\frac a2\right)^2-\left(\frac{\sqrt2}{4}\right)^2}=\frac12\sqrt{a^2-\frac12}
\]
所以可取
\(E=\left(0,\frac{a+\sqrt{a^2-\frac12}}2\right)\)，\(F=\left(0,\frac{a-\sqrt{a^2-\frac12}}2\right)\)
且 \(|PE|+|PF|=a\).
\end{solution}

\begin{problem}
设 \(a_0\) 为常数，且 \(a_n = 3^{n-1} - 2a_{n-1}\)（\(n \in \N^*\)）.

\begin{enumerate}
\item 证明对任意 \(n \geqslant 1\)，\(a_n = \frac{1}{5} \left[3^n + (-1)^{n-1} \cdot 2^n\right] + (-1)^n \cdot 2^n a_0\)；
\item 假设对任意 \(n \geqslant 1\) 有 \(a_n > a_{n-1}\)，求 \(a_0\) 的取值范围.
\end{enumerate}
\end{problem}

\begin{solution}
\begin{enumerate}
\item 先验证初始情形. 由递推式
\[
a_1=1-2a_0
\]
可得
\[
a_1-a_0=1-3a_0
\]
设对某个 \(n\geqslant1\) 已有
\[
a_n=\frac15\left[3^n+(-1)^{n-1}2^n\right]+(-1)^n2^na_0
\]
则
\[
a_{n+1}=3^n-2a_n=\frac15\left[3^{n+1}+(-1)^n2^{n+1}\right]+(-1)^{n+1}2^{n+1}a_0
\]
这正是把 \(n\) 换成 \(n+1\) 后的公式，结合 \(n=1\) 时公式成立，归纳得
\[
a_n=\frac15\left[3^n+(-1)^{n-1}2^n\right]+(-1)^n2^na_0
\]

\item 由上式计算相邻两项之差，得
\[
a_n-a_{n-1}=\frac{2\cdot3^{n-1}+3(-1)^{n-1}2^{n-1}(1-5a_0)}5
\]
取 \(n=1\)，\(2\)，分别得到
\(a_1-a_0=1-3a_0>0\)，\(a_2-a_1=6a_0>0\)
所以必须有 \(0<a_0<\dfrac13\).

反之，设 \(0<a_0<\dfrac13\). 当 \(n\) 为奇数且 \(n\geqslant3\) 时，由 \(1-5a_0>-\dfrac23\)，有
\[
5(a_n-a_{n-1})>2\cdot3^{n-1}-2\cdot2^{n-1}>0
\]
当 \(n=1\) 时，前面已经有 \(a_1-a_0>0\). 当 \(n\) 为偶数且 \(a_0\geqslant\dfrac15\) 时，\(1-5a_0\leqslant0\)，所以
\[
5(a_n-a_{n-1})=2\cdot3^{n-1}-3\cdot2^{n-1}(1-5a_0)\geqslant2\cdot3^{n-1}>0
\]
若 \(0<a_0<\dfrac15\)，则
\[
5(a_n-a_{n-1})>2\cdot3^{n-1}-3\cdot2^{n-1}\geqslant0
\]
且在可能取等的 \(n=2\) 时因 \(a_0>0\) 实际为正. 因此对任意 \(n\geqslant1\) 都有 \(a_n>a_{n-1}\).

所以 \(a_0\) 的取值范围为
\[
\left(0,\frac13\right)
\]
\end{enumerate}
\end{solution}
